\section{System: Implementation, Tests and Limitations}\label{app:system}

\textbf{Code.} \texttt{src/forgery/pipeline.py} (\texttt{detect()}), \texttt{src/forgery/explain.py} (explanations,
input preparation and reports; no interface code), \texttt{detect.py} (command line) and \texttt{app/} (Streamlit
interface, started with \texttt{streamlit run app/Home.py}). The PDF report holds the verdict, overlay, evidence chart
and caveats on the first page and the eight evidence maps on the second. The command-line tool and both analysis pages
use the same input preparation. The detector's loader converts every image with PIL's \texttt{convert("RGB")}; this
is correct for CASIA, where all 12,614 images are 8-bit RGB without EXIF rotation, but it would turn a 16-bit greyscale
image white. A changed image is saved losslessly as PNG; under B and R the detector works on decoded pixels, so this
does not affect the analysis. The explorer shows only the 10,722 training and validation images, and since the
deployed models were refit on all development images, their results on CASIA examples are optimistic.

\textbf{Input preparation.} Every user image is fully decoded, so truncated files fail early; decompression bombs and
images outside 64--4,000~px per side are refused; the image is rotated by its EXIF orientation, and 16-bit or float
images are rescaled to 8~bits.

\textbf{Performance.} Table~\ref{tab:sys-perf} gives run times on an Apple M4 laptop. Models are loaded once and
cached.

\begin{table}[!ht]
\centering
\caption{Run time on an Apple M4 laptop (largest allowed image: $4{,}000\times2{,}667$~px).}
\label{tab:sys-perf}
\footnotesize
\begin{tabular}{lcc}
\toprule
Step & CASIA-sized & Largest (R / B) \\
\midrule
Detection, all eight modules & $\approx$0.2~s & 4.3~s / 7.1~s \\
SHAP explanation & $\approx$0.5~s & $\approx$0.5~s \\
PDF report & $\approx$0.3~s & 2.0~s / 3.3~s \\
\bottomrule
\end{tabular}
\end{table}

\textbf{Interface details.} The analysis page accepts JPEG, PNG, TIFF, BMP and WebP (at most 4,000~px per side) and, for
CASIA examples, also shows the ground truth. The batch page shows a progress bar, a table of verdicts, probabilities,
region sizes and the strongest evidence, a gallery of images judged tampered, and CSV and ZIP (overlay) downloads. The
dataset explorer shows ground-truth masks and metadata (format, estimated JPEG quality, IJG tables, tampered area,
host and donor id, split group). The results dashboard covers leakage, classification, localisation, the test set,
robustness and MICC-F220, and the CNN comparison; it shows saved figures made from test images (the missed forgeries
and false alarms of Section~\ref{sec:test} and the synthetic forgeries of Section~\ref{sec:robustness}), which are
outputs of the one-time evaluation, not a way to browse the test set. Uploads are stored once per content hash, so a
rerun (a protocol switch back or a download click) reuses the cached analysis and PDF, also for rotated or 16-bit
photos, whose prepared copy is written once; batch results are kept in the session until the input or the protocol
changes; stored uploads are deleted after seven days.

\textbf{Tests.} \texttt{tests/test\_app.py} runs every page headlessly with Streamlit's AppTest framework and checks
that no page raises an exception; the analysis page produces a verdict for an example image; the batch page processes
a folder containing a corrupt file and reports exactly that file as an error; a batch in which every file fails still
renders and hidden files are skipped; batch results survive a rerun; and an empty explorer filter does not blank the
other tabs. \texttt{tests/test\_explain.py} checks that 16-bit images are rescaled, not clipped; EXIF rotation is
applied; truncated, non-image, too-small and too-large files are refused; and ordinary files pass through unchanged.
All pages were also exercised in headless Chrome, including a batch run and the example analysis; the screenshots come
from that run.

\textbf{Limitations.} (i)~Local use only: nothing is uploaded to a server, but the folder mode reads any folder the user
names, so the configuration binds the server to \texttt{localhost}. (ii)~An EXIF-rotated photo is analysed upright;
several modules work on fixed block grids, so the same photo rotated by 90$^\circ$ can receive a different score.
(iii)~The suspected region is shown for uncertain images too, for inspection, matching the command-line tool, which
removes the region only for images judged authentic.
